\section{One-step governor derivation and guarantee boundary}
\label{app:governor-derivation}

\subsection{Affine dynamics and scalar objective}

% CLAIM: C09
Consider one axis with current state
\(x=[p,v,a]^{\mathsf T}\), step \(\Delta=\DT\), and constant jerk \(j\).
Exact integration yields
\begin{equation}
  x^+(j)=b+gj,
  \qquad
  b=
  \begin{bmatrix}
    p+v\Delta+\tfrac12a\Delta^2\\
    v+a\Delta\\
    a
  \end{bmatrix},
  \qquad
  g=
  \begin{bmatrix}
    \Delta^3/6\\ \Delta^2/2\\ \Delta
  \end{bmatrix}.
  \label{eq:appendix-affine-dynamics}
\end{equation}
For objective \(\rho=[\rho_p,\rho_v,\rho_a]^{\mathsf T}\), define
\(\tilde w_c=w_c/s_c^2\).  Expanding
\eqref{eq:governor-objective} gives
\begin{equation}
  \mathcal L(j)=Qj^2+2Lj+C,
  \label{eq:appendix-quadratic}
\end{equation}
where
\begin{align}
  Q &=
  \sum_c \tilde w_c g_c^2+
  \frac{\lambda_j+\lambda_{\Delta j}}{\JerkLimitSymbol^2},
  \label{eq:appendix-Q}\\
  L &=
  \sum_c \tilde w_c g_c(b_c-\rho_c)
  -\frac{\lambda_{\Delta j}j_{k-1}}{\JerkLimitSymbol^2}.
  \label{eq:appendix-L}
\end{align}
The implementation forms the preferred unconstrained value
\(j^{\mathrm{unc}}=-L/\max(Q,\epsilon_Q)\), with a small positive numerical
floor \(\epsilon_Q\), before applying analytic interval constraints.

\subsection{Continuous segment constraints}

Let \(V=\VelocityLimitSymbol\), \(A=\AccelerationLimitSymbol\),
\(J=\JerkLimitSymbol\), and suppose the current state is finite and satisfies
\(\lvert v\rvert\leq V\) and \(\lvert a\rvert\leq A\).  Jerk, terminal
acceleration, and terminal velocity first give the interval
\begin{equation}
  \mathcal J_{\mathrm{end}}(x)=
  [-J,J]
  \cap
  \left[\frac{-A-a}{\Delta},\frac{A-a}{\Delta}\right]
  \cap
  \left[
    \frac{2(-V-v-a\Delta)}{\Delta^2},
    \frac{2(V-v-a\Delta)}{\Delta^2}
  \right].
  \label{eq:endpoint-jerk-interval}
\end{equation}

% CLAIM: C09
Endpoint velocity checks are insufficient when acceleration changes sign
inside the interval.  Because
\(v(\tau)=v+a\tau+\tfrac12j\tau^2\), the only possible interior extremum is
\(\tau^\star=-a/j\).  For \(a>0\) and \(j<-a/\Delta\), this is an interior
maximum and requires
\begin{equation}
  v-\frac{a^2}{2j}\leq V.
  \label{eq:positive-acceleration-extremum}
\end{equation}
For \(a<0\) and \(j>-a/\Delta\), it is an interior minimum and requires
\begin{equation}
  v-\frac{a^2}{2j}\geq -V.
  \label{eq:negative-acceleration-extremum}
\end{equation}
Splitting at \(j=-a/\Delta\) turns these rational conditions into interval
thresholds.  Their intersection with
\(\mathcal J_{\mathrm{end}}(x)\) is the possibly disconnected set
\(\mathcal J_{\mathrm{seg}}(x)\) of jerks satisfying continuous V/A/J bounds.
No time or jerk grid is used to construct this set.

\subsection{Stopping viability and next-step existence}

Define the point-admissible set
\begin{equation}
  \mathcal A =
  \{[p,v,a]^{\mathsf T}:
  p,v,a\ \text{finite},\ |v|\leq V,\ |a|\leq A\}
  \label{eq:point-admissible-set}
\end{equation}
and the directional stopping-viability set
\begin{equation}
  \mathcal V =
  \left\{x\in\mathcal A:
  \begin{array}{ll}
    v+a^2/(2J)\leq V,&a>0,\\
    v-a^2/(2J)\geq -V,&a<0
  \end{array}
  \right\}.
  \label{eq:viability-set}
\end{equation}
For \(a=0\), point admissibility is sufficient for membership in
\(\mathcal V\).

% CLAIM: C09
The analytic one-step interval is
\begin{equation}
  \mathcal J_1(x)=
  \{j\in\mathcal J_{\mathrm{seg}}(x):x^+(j)\in\mathcal V\},
  \qquad x\in\mathcal V.
  \label{eq:one-step-viable-jerks}
\end{equation}
The sign of \(a^+(j)=a+j\Delta\) splits the terminal envelope into two
quadratic inequalities in \(j\).  Analytic roots are intersected with
\(\mathcal J_{\mathrm{seg}}(x)\), producing a union of closed intervals.
The implementation additionally verifies
\begin{equation}
  \mathcal J_1\bigl(x^+(j)\bigr)\neq\varnothing,
  \label{eq:next-step-existence}
\end{equation}
so the accepted endpoint has at least one viable action at the next tick.  It
is useful to denote the postchecked subset by
\begin{equation}
  \mathcal J_{\mathrm{next}}(x)=
  \{j\in\mathcal J_1(x):
  \mathcal J_1(x^+(j))\neq\varnothing\}.
  \label{eq:next-preserving-jerks}
\end{equation}

\subsection{Selection and recovery procedure}

For intervals
\(\mathcal J_1(x)=\bigcup_r[\ell_r,u_r]\), the normal procedure is:
\begin{enumerate}
  \item clip \(j^{\mathrm{unc}}\) to every interval and retain the closest
  candidate;
  \item integrate the candidate exactly and verify segment legality, terminal
  stopping viability, and \eqref{eq:next-step-existence};
  \item if that candidate fails a numerical boundary postcheck, evaluate a
  deterministic recovery set containing the clipped zero-terminal-acceleration
  action, interval boundaries, neighbouring floating-point values, and
  interval midpoints;
  \item among verified recovery candidates, choose the lowest value of
  \(\mathcal L\); if none exists, report an invariant violation rather than
  committing an unverified state.
\end{enumerate}
The recovery set is a boundary-robust implementation mechanism, not a sampled
feasibility oracle.  Because it is finite, the recovery branch must not be
described as a global minimizer over
\(\mathcal J_{\mathrm{next}}(x)\).

For a nonfinite objective and a viable current state, the implementation
instead seeks a verified action near \(j=-a/\Delta\), favouring reduced
terminal acceleration, then reduced terminal velocity and jerk.  If the
current state is outside the normal viability conditions, a bounded
best-effort emergency jerk opposes the approached velocity boundary or removes
acceleration.  This emergency action remains dynamically integrated, but its
endpoint is not assigned the normal guarantee.

\subsection{Conditional guarantee and exclusions}

% CLAIM: C09
For any returned non-emergency action whose postchecks succeed, the
implementation verifies all of the following under the triple-integrator
model:
\begin{enumerate}
  \item the committed endpoint equals \(x^+(j)\);
  \item velocity, acceleration, and jerk limits hold continuously over the
  command interval;
  \item the endpoint lies in \(\mathcal V\); and
  \item at least one subsequent jerk belongs to
  \(\mathcal J_1(x^+(j))\).
\end{enumerate}
For multiple axes with uncoupled componentwise limits, the construction uses
the Cartesian product of the per-axis verified sets and a common \(\DT\).
This does not add coupled path, synchronization-optimality, torque, or
collision constraints.

The statement is conditional on positive known limits and control period,
finite current state, a viable selected current state for the normal path,
exact execution of the selected constant jerk, the declared current-state
feedback policy, and no unmodelled disturbance during the interval.
Next-step existence is a one-step property; it is not a theorem of recursive
feasibility for arbitrary future corrections or disturbances.  The objective
is an implementation tracking preference, not a proof of global trajectory
optimality, and the verified command-state envelope is not a functional-safety
certificate.

